\section{Experimental Setup}

\subsection{Compared Methods}

The experiment compares three methods:

\begin{enumerate}
    \item \textbf{Lexicon}: the weighted NRC Hashtag Emotion Lexicon scorer.
    \item \textbf{NB}: the weakly supervised Naive Bayes model trained from lexicon entries.
    \item \textbf{Hybrid}: score-level combination of lexicon and Naive Bayes scores.
\end{enumerate}

The same preprocessing pipeline and emotion label set are used across all methods.

\subsection{Dataset}

The evaluation dataset contains 24 labeled texts. It is balanced across the eight emotion labels, with three examples per emotion. It also includes a \texttt{domain} column with three values: \texttt{social}, \texttt{news}, and \texttt{blog}. Each domain contains 8 examples.

\begin{table}[h]
\centering
\small
\begin{tabular}{lrl}
\toprule
\textbf{Property} & \textbf{Value} & \textbf{Description} \\
\midrule
Rows & 24 & total examples \\
Emotion labels & 8 & anger, anticipation, disgust, fear, joy, sadness, surprise, trust \\
Examples per emotion & 3 & balanced label distribution \\
Domains & 3 & social, news, blog \\
Examples per domain & 8 & balanced domain distribution \\
\bottomrule
\end{tabular}
\caption{Evaluation dataset summary.}
\label{tab:dataset}
\end{table}

\subsection{Metrics}

The main metrics are accuracy and macro-F1. Accuracy measures the fraction of examples where the predicted dominant emotion matches the gold label. Macro-F1 is the average of the F1 scores computed independently for each emotion. Macro-F1 is useful for multi-class emotion analysis because it treats all emotions equally.

The experiment also writes confusion matrices and per-domain metrics to support qualitative analysis.
